\pdfbookmark{Общая характеристика работы}{characteristic}             % Закладка pdf
\section*{Общая характеристика работы}

\newcommand{\actuality}{\pdfbookmark[1]{Актуальность}{actuality}\underline{\textbf{\actualityTXT}}}
\newcommand{\progress}{\pdfbookmark[1]{Разработанность темы}{progress}\underline{\textbf{\progressTXT}}}
\newcommand{\aim}{\pdfbookmark[1]{Цели}{aim}\underline{{\textbf\aimTXT}}}
\newcommand{\tasks}{\pdfbookmark[1]{Задачи}{tasks}\underline{\textbf{\tasksTXT}}}
\newcommand{\aimtasks}{\pdfbookmark[1]{Цели и задачи}{aimtasks}\aimtasksTXT}
\newcommand{\object}{\pdfbookmark[1]{Объект исследования}{object}\underline{\textbf{\objectTXT}}}
\newcommand{\subject}{\pdfbookmark[1]{Предмет исследования}{subject}\underline{\textbf{\subjectTXT}}}
\newcommand{\novelty}{\pdfbookmark[1]{Научная новизна}{novelty}\underline{\textbf{\noveltyTXT}}}
\newcommand{\tinfluence}{\pdfbookmark[1]{Теоретическая значимость}{tinfluence}\underline{\textbf{\tinfluenceTXT}}}
\newcommand{\pinfluence}{\pdfbookmark[1]{Практическая значимость}{pinfluence}\underline{\textbf{\pinfluenceTXT}}}
\newcommand{\methods}{\pdfbookmark[1]{Методология и методы исследования}{methods}\underline{\textbf{\methodsTXT}}}
\newcommand{\defpositions}{\pdfbookmark[1]{Положения, выносимые на защиту}{defpositions}\underline{\textbf{\defpositionsTXT}}}
\newcommand{\passport}{\pdfbookmark[1]{Соответствие паспорту специальности}{passport}\underline{\textbf{\passportTXT}}}
\newcommand{\reliability}{\pdfbookmark[1]{Достоверность}{reliability}\underline{\textbf{\reliabilityTXT}}}
\newcommand{\probation}{\pdfbookmark[1]{Апробация}{probation}\underline{\textbf{\probationTXT}}}
\newcommand{\contribution}{\pdfbookmark[1]{Личный вклад}{contribution}\underline{\textbf{\contributionTXT}}}
\newcommand{\publications}{\pdfbookmark[1]{Публикации}{publications}\underline{\textbf{\publicationsTXT}}}

\providecommand{\extcite}[1]{\ifsynopsis\else~\cite{#1}\fi}% внешние ссылки только в диссертации
{\actuality} В статистической теории обучения обоснована принципиальная связь между сложностью семейства моделей, объемом обучающей выборки и способностью к обобщению\extcite{vapnik1995nature}. Для надежного выбора модели требуется достаточное количество данных. Подавляющее большинство современных моделей машинного обучения являются параметрическими: они задаются конечным набором параметров, определяемых минимизацией эмпирической функции потерь на выборке. Поскольку функция потерь вычисляется по ограниченным данным и зависит от параметров модели, возникает фундаментальный вопрос о взаимосвязи объема выборки, структуры модели и свойств множества оптимальных параметров, а также о \textbf{минимальном объеме данных}, необходимом для успешного обучения\extcite{kaplan2020scaling,hoffmann2022}.

Эмпирические \textbf{законы масштабирования} показывают, что качество моделей закономерно зависит от размера данных и модели: в языковых моделях установлены степенные зависимости потерь от объема данных и числа параметров\extcite{kaplan2020scaling}, а для вычислительно-оптимального обучения предложено согласованное масштабирование размера модели и объема данных\extcite{hoffmann2022}. Эти результаты опираются на масштабные эксперименты и не имеют пока единого теоретического обоснования.

Естественным объектом для такого анализа является \textbf{ландшафт функции потерь} в пространстве параметров. Его визуализация и анализ структуры минимумов\extcite{li2018visualizing}, изучение локальной геометрии и связей с обобщающей способностью\extcite{fort2019large}, а также доказательство связности множества решений\extcite{draxler2018essentially} показывают, что геометрия поверхности функции потерь определяет как выбор решения методами градиентной оптимизации\extcite{garipov2018loss,chaudhari2019}, так и устойчивость к возмущениям данных. Спектральные свойства \textbf{матрицы Гессе} в окрестности минимума характеризуют кривизну оптимизационной поверхности и используются при анализе сходимости и обобщения\extcite{papyan2019}. Вместе с тем остается открытой задача систематического описания того, как ландшафт функции потерь стабилизируется при увеличении объема выборки и как связать эту стабилизацию с законами масштабирования и достаточным размером данных.

\ifsynopsis\textit{Таким образом, актуальной задачей является разработка теоретических методов, связывающих архитектуру параметрической модели, объем обучающей выборки и локальную геометрию множества оптимальных параметров, задаваемую ландшафтом функции потерь.}\else{Таким образом, актуальной задачей является разработка теоретических методов, связывающих архитектуру параметрической модели, объем обучающей выборки и локальную геометрию множества оптимальных параметров, задаваемую ландшафтом функции потерь.}\fi

\ifsynopsis
{\progress} Задача определения достаточного размера выборки традиционно решается методами математической статистики, где размер выборки выбирается так, чтобы критерий проверки гипотез имел заданную мощность при фиксированной вероятности ошибки первого рода, методами байесовского анализа, основанными на апостериорной неопределенности и близости апостериорных распределений, и методами статистической теории обучения, связывающими объем выборки со сложностью класса моделей. Для современных нейронных сетей эти подходы либо требуют явных гипотез о распределении параметров, либо дают сильно завышенные оценки. Ландшафт функции потерь и спектр матрицы Гессе активно исследуются в связи с оптимизацией и обобщающей способностью нейронных сетей, однако связь локальной геометрии оптимизационной поверхности с достаточным размером выборки и законами масштабирования остается открытым вопросом.
\else
\input{intro_review}
\fi

{\aim} данной работы является разработка ландшафтного метода оценки масштабируемости параметрических моделей машинного обучения, основанного на локальной геометрии ландшафта функции потерь в окрестности точки минимума и на стабилизации распределения оптимальных параметров. Для достижения цели были поставлены и решены следующие {\tasks}:

\begin{enumerate}[beginpenalty=10000]
    \item Исследовать спектральные свойства матриц Гессе нейронных сетей и получить архитектурно-зависимые верхние оценки спектральных норм для полносвязных и сверточных архитектур.
    \item Разработать критерии стабилизации оптимизационной поверхности при увеличении объема обучающей выборки.
    \item Оценить масштабируемость параметрических моделей машинного обучения на основе локальной геометрии ландшафта функции потерь, связав ее с архитектурой модели и объемом обучающей выборки.
    \item Построить критерии достаточного размера выборки для параметрических моделей на основе стабилизации ландшафта функции потерь.
    \item Исследовать вероятностные линейные модели как частный случай параметрических моделей и построить для них критерии достаточного размера выборки на основе сходимости распределения оптимальных параметров.
\end{enumerate}

{\object} настоящей работы являются параметрические модели машинного обучения и возникающие при их обучении оптимизационные задачи на конечных обучающих выборках.

{\subject} данной работы являются локальные свойства оптимизационной поверхности функции потерь параметрических моделей машинного обучения, спектральные свойства матрицы Гессе, стабилизация множества оптимальных параметров при увеличении объема обучающей выборки и критерии достаточного размера выборки, основанные на этой стабилизации.

{\novelty}
\begin{enumerate}[beginpenalty=10000]
    \item Получены верхние оценки спектральных норм гаусс-ньютоновской компоненты матрицы Гессе для полносвязных и сверточных архитектур нейронных сетей на основе ее блочной факторизации в классе матрично-представимых моделей.
    \item Предложены критерии стабилизации оптимизационной поверхности. Порядок их убывания следует из усреднения эмпирического риска по объектам выборки; локальная квадратичная аппроксимация и спектр матрицы Гессе оценивают константу в этой оценке.
    \item Предложена схема законов масштабирования, в которой ошибка оценивания задается отношением шума градиентов к кривизне. Для линейной модели с неполным набором признаков схема выполнена точно.
    \item Получены критерии достаточного размера выборки для параметрических моделей на основе стабилизации ландшафта функции потерь.
    \item Для вероятностных линейных моделей построены критерии достаточного размера выборки, основанные на сходимости апостериорного распределения параметров и стабилизации вероятностных характеристик решения.
\end{enumerate}

{\tinfluence} полученных результатов состоит в том, что предложено обоснование связи между объемом обучающей выборки и стабилизацией оптимизационной поверхности (для линейных моделей "--- через правдоподобие и апостериорные распределения, для нейронных сетей "--- через спектр матрицы Гессе и сходимость ландшафта функции потерь), дополняющее известные эмпирические законы масштабирования строгими оценками и единой постановкой задачи.

{\pinfluence} настоящей работы заключается в возможности применения разработанных методов для обоснованного определения достаточного объема данных при планировании экспериментов и сборе выборок для линейных и нейросетевых моделей, а также для формулировки критериев остановки сбора данных и интерпретации достаточности выборки в терминах геометрии поверхности функции потерь.

{\methods} Для достижения поставленных
в диссертационной работе целей используются:
\begin{enumerate}[beginpenalty=10000]
    \item Методы линейной алгебры, математического и функционального анализа: матричное дифференцирование, анализ функций многих переменных, спектральный анализ матриц (в том числе матрицы Гессе).
    \item Методы теории вероятностей и математической статистики: байесовский вывод, анализ сходимости последовательностей случайных величин и распределений, элементы теории оптимизации.
    \item Методы вычислительного эксперимента и анализа результатов: бутстрап (англ. bootstrap), метод Монте"--~Карло, воспроизводимая постановка экспериментов.
\end{enumerate}

{\defpositions}
\begin{enumerate}[beginpenalty=10000]
    \item Теоремы о верхних оценках спектральной нормы гаусс-ньютоновской компоненты матрицы Гессе для полносвязных и сверточных архитектур нейронных сетей.
    \item Критерии стабилизации ландшафта функции потерь: порядок $\mathcal O(k^{-1})$ и $\mathcal O(k^{-2})$ следует из усреднения по объектам, а оценки констант получены по локальной квадратичной аппроксимации в одной точке, в полном пространстве параметров и в подпространстве главных собственных направлений матрицы Гессе.
    \item Схема законов масштабирования, в которой ошибка оценивания задается отношением шума градиентов к кривизне, и ее точный вид для линейной модели с неполным набором признаков.
    \item Критерии достаточного размера выборки для параметрических моделей, основанные на стабилизации ландшафта функции потерь.
    \item Критерии достаточного размера выборки для вероятностных линейных моделей, основанные на стабилизации распределения оптимальных параметров при увеличении объема данных.
\end{enumerate}

{\passport} Диссертационная работа соответствует паспорту научной специальности 1.2.1 "--- <<Искусственный интеллект и машинное обучение>>, а именно в наибольшей степени следующим направлениям исследований: п.~1 <<Естественно-научные основы и методы искусственного интеллекта>>; п.~2 <<Исследования в области оценки качества и эффективности алгоритмических решений для систем искусственного интеллекта и машинного обучения>>; п.~12 <<Исследования в области надежности, устойчивости и переобучения систем класса ИИ>>; п.~15 <<Математические исследования в области статистики, логики, алгебры, топологии, анализа функций и других областях, ориентированные на решение задач искусственного интеллекта и машинного обучения>>; п.~16 <<Исследования в области специальных методов оптимизации, проблем сложности и снижения размерности>>; п.~17 <<Исследования в области многослойных алгоритмических конструкций, в том числе многослойных нейросетей>>.

{\reliability} полученных результатов обеспечивается корректностью
применения апробированного в научной практике математического аппарата теории вероятностей, математической статистики и других смежных теоретических областей, а также экспериментальной проверкой разработанных методов на большом количестве модельных и практических задач машинного обучения. Полученные теоретические результаты обосновываются математически строгими доказательствами, а для проведенных вычислительных экспериментов даются детальные описания, обеспечивающие их воспроизводимость. Результаты опубликованы в рецензируемых научных изданиях и в трудах рецензируемых российских и международных конференций по машинному обучению и искусственному интеллекту.

{\probation} Основные результаты работы докладывались и обсуждались на следующих научных конференциях:
\begin{enumerate}
    \item 16-я Международная конференция
    <<Интеллектуализация обработки информации>> (ИОИ-2026),
    Витебск, 28~сентября~-- 2~октября 2026~г.
    \item Международная конференция <<Иванниковские чтения>>,
    Этномир, 28--30~мая 2026~г.
    \item 68-я Всероссийская научная конференция МФТИ,
    Москва, 30~марта~-- 4~апреля 2026~г.
    \item Открытая конференция ИСП РАН, Москва, 9--10 декабря 2025~г.
    \item 22-я Всероссийская конференция с международным участием
    <<Математические методы распознавания образов (ММРО-2025)>>,
    Муром, 22--26 сентября 2025~г.
    \item 67-я Всероссийская научная конференция МФТИ,
    Москва, 31 марта -- 5 апреля 2025~г.
    \item Открытая конференция ИСП РАН, Москва, 11--12 декабря 2024~г.
    \item 66-я Всероссийская научная конференция МФТИ,
    Москва, 1--6 апреля 2024~г.
\end{enumerate}


{\contribution} Все приведенные результаты, кроме отдельно оговоренных случаев, получены диссертантом лично при научном руководстве к.ф.-м.н.~А.\,В.~Грабового.

\ifnumequal{\value{bibliosel}}{0}
{%%% Встроенная реализация с загрузкой файла через движок bibtex8. (При желании, внутри можно использовать обычные ссылки, наподобие `\cite{vakbib1,vakbib2}`).
    {\publications} Основные результаты по теме диссертации изложены
    в~XX~печатных изданиях,
    X из которых изданы в журналах, рекомендованных ВАК,
    X "--- в тезисах докладов.
}%
{%%% Реализация пакетом biblatex через движок biber
    \begin{refsection}[bl-author, bl-registered]
        % Это refsection=1.
        % Процитированные здесь работы:
        %  * подсчитываются, для автоматического составления фразы "Основные результаты ..."
        %  * попадают в авторскую библиографию, при usefootcite==0 и стиле `\insertbiblioauthor` или `\insertbiblioauthorgrouped`
        %  * нумеруются там в зависимости от порядка команд `\printbibliography` в этом разделе.
        %  * при использовании `\insertbiblioauthorgrouped`, порядок команд `\printbibliography` в нем должен быть тем же (см. biblio/biblatex.tex)
        %
        % Невидимый библиографический список для подсчета количества публикаций:
        \phantom{\printbibliography[heading=nobibheading, section=1, env=countauthorvak,          keyword=biblioauthorvak]%
        \printbibliography[heading=nobibheading, section=1, env=countauthorwos,          keyword=biblioauthorwos]%
        \printbibliography[heading=nobibheading, section=1, env=countauthorscopus,       keyword=biblioauthorscopus]%
        \printbibliography[heading=nobibheading, section=1, env=countauthorconf,         keyword=biblioauthorconf]%
        \printbibliography[heading=nobibheading, section=1, env=countauthorother,        keyword=biblioauthorother]%
        \printbibliography[heading=nobibheading, section=1, env=countauthor,             keyword=biblioauthor]%
        \printbibliography[heading=nobibheading, section=1, env=countauthorvakscopuswos, filter=vakscopuswos]%
        \printbibliography[heading=nobibheading, section=1, env=countauthorscopuswos,    filter=scopuswos]}%
        %
        \nocite{*}%
        %
        % {\publications} Основные результаты по теме диссертации изложены в~\arabic{citeauthor}~печатных изданиях~\cite{},
        % \arabic{citeauthorvak} из которых изданы в журналах, рекомендованных ВАК%
        % \ifnum \value{citeauthorscopuswos}>0%
        %     , \arabic{citeauthorscopuswos} "--- в~периодических научных журналах, индексируемых Web of~Science и Scopus%
        % \fi%
        % \ifnum \value{citeauthorconf}>0%
        %     , \arabic{citeauthorconf} "--- в~тезисах докладов.
        % \else%
        %     .
        % \fi%
        % \ifnum \value{citeregistered}=1%
        %     \ifnum \value{citeauthorpatent}=1%
        %         Зарегистрирован \arabic{citeauthorpatent} патент.
        %     \fi%
        %     \ifnum \value{citeauthorprogram}=1%
        %         Зарегистрирована \arabic{citeauthorprogram} программа для ЭВМ.
        %     \fi%
        % \fi%
        % \ifnum \value{citeregistered}>1%
        %     Зарегистрированы\ %
        %     \ifnum \value{citeauthorpatent}>0%
        %     \formbytotal{citeauthorpatent}{патент}{}{а}{}%
        %     \ifnum \value{citeauthorprogram}=0 . \else \ и~\fi%
        %     \fi%
        %     \ifnum \value{citeauthorprogram}>0%
        %     \formbytotal{citeauthorprogram}{программ}{а}{ы}{} для ЭВМ.
        %     \fi%
        % \fi%
        {\publications} \ifsynopsis\textit{Список публикаций приведен в конце автореферата. }\else\textit{Работы автора приведены в общем списке литературы. }\fi Основные результаты по теме диссертации изложены в~\arabic{citeauthor}~научных работах, из которых~\arabic{citeauthorscopuswos} "--- в~рецензируемых периодических научных журналах, а остальные~\arabic{citeauthorconf} "--- в~сборниках трудов конференций и тезисах докладов.
        % К публикациям, в которых излагаются основные научные результаты диссертации на соискание ученой
        % степени, в рецензируемых изданиях приравниваются патенты на изобретения, патенты (свидетельства) на
        % полезную модель, патенты на промышленный образец, патенты на селекционные достижения, свидетельства
        % на программу для электронных вычислительных машин, базу данных, топологию интегральных микросхем,
        % зарегистрированные в установленном порядке.(в ред. Постановления Правительства РФ от 21.04.2016 N 335)
    \end{refsection}%
    \begin{refsection}[bl-author, bl-registered]
        % Это refsection=2.
        % Процитированные здесь работы:
        %  * попадают в авторскую библиографию, при usefootcite==0 и стиле `\insertbiblioauthorimportant`.
        %  * ни на что не влияют в противном случае
        \nocite{kiselev2026ispras}
        \nocite{kiselev2025ssdpdp}
        \nocite{kiselev2025ssdlb}
        \nocite{kiselev2024unraveling}
        \nocite{kiselev2026conf68mipt}
        \nocite{kiselev2025mmro}
        \nocite{kiselev2025conf67mipt}
        \nocite{meshkov2024convnets}
        \nocite{kiselev2024conf66mipt}
    \end{refsection}%
        %
        % Все, что вне этих двух refsection, это refsection=0,
        %  * для диссертации - это нормальные ссылки, попадающие в обычную библиографию
        %  * для автореферата:
        %     * при usefootcite==0, ссылка корректно сработает только для источника из `external.bib`. Для своих работ "--- напечатает "[0]" (и даже Warning не вылезет).
        %     * при usefootcite==1, ссылка сработает нормально. В авторской библиографии будут только процитированные в refsection=0 работы.
} % Характеристика работы по структуре во введении и в автореферате не отличается (ГОСТ Р 7.0.11, пункты 5.3.1 и 9.2.1), потому её загружаем из одного и того же внешнего файла, предварительно задав форму выделения некоторым параметрам

%Диссертационная работа была выполнена при поддержке грантов \dots

\InputIfFileExists{common/dissvolume.tex}{}{%
    \PackageWarning{synopsis}{Нет common/dissvolume.tex: сначала соберите диссертацию}%
}
\underline{\textbf{Объем и структура работы.}} Диссертация состоит из введения,
\formbytotal{totalchapter}{глав}{ы}{}{} и заключения. % родительный: 1 главы, 2 глав, 5 глав
Полный объём диссертации составляет
\formbytotal{TotPages}{страниц}{у}{ы}{}, включая
\formbytotal{totalcount@figure}{рисун}{ок}{ка}{ков} и
\formbytotal{totalcount@table}{таблиц}{у}{ы}{}.
Список литературы содержит
\formbytotal{citeexternal}{наименован}{ие}{ия}{ий}.

\pdfbookmark{Содержание работы}{description}                          % Закладка pdf
\section*{Содержание работы}
Во \underline{\textbf{введении}} обосновывается актуальность
исследований, проводимых в~рамках данной диссертационной работы,
приводится обзор научной литературы по~изучаемой проблеме,
формулируется цель, ставятся задачи работы, излагается научная новизна
и практическая значимость представляемой работы, а также вводятся основные понятия. Объектом называется пара $(\mathbf{x}, y)$, где $\mathbf{x} \in \mathbb{X} \subseteq \mathbb{R}^n$ "--- вектор признакового описания, $y \in \mathbb{Y}$ "--- целевая переменная. В задаче регрессии $\mathbb{Y} = \mathbb{R}$, в задаче классификации $\mathbb{Y} = \{1, \ldots, K\}$. Выборка размера~$m$ задается множеством $\mathfrak{D}_m = \{(\mathbf{x}_i, y_i)\}_{i=1}^m$ из $m$ независимых реализаций пар~$(\mathbf{x}, y)$, распределенных согласно неизвестному распределению $p(\mathbf{x}, y)$.

Параметрической моделью называется семейство функций $f_{\mathbf{w}}: \mathbb{X} \to \mathbb{Y}$, задаваемое вектором параметров $\mathbf{w} \in \mathbb{W} \subset \mathbb{R}^N$. Выбор модели из этого семейства проводится с помощью эмпирической функции потерь, вычисляемой на заданной выборке:
\begin{equation}
\label{eq:synopsis-empirical-loss}
\mathcal{L}_m(\mathbf{w}) = \frac{1}{m} \sum_{i=1}^m \ell(f_{\mathbf{w}}(\mathbf{x}_i), y_i),
\end{equation}
где $\ell: \mathbb{Y} \times \mathbb{Y} \to \mathbb{R}^+$ "--- функция потерь на одном объекте (например, квадратичная или кросс-энтропийная). Обучение модели таким образом есть решение оптимизационной задачи
\begin{equation}
    \label{eq:synopsis-min-problem}
    \mathbf{w}_m^* \in \argmin_{\mathbf{w} \in \mathbb{W}} \mathcal{L}_m(\mathbf{w}).
\end{equation}

В качестве ключевого объекта исследования предложено использовать оптимизационную поверхность функции потерь и ее локальную геометрию в окрестности точки минимума. Эта локальная геометрия, а не глобальная картина поверхности, составляет основу ландшафтного метода оценки масштабируемости: стабилизация ландшафта при увеличении объема обучающей выборки дает критерии достаточного размера выборки и схему законов масштабирования.


\underline{\textbf{В первой главе}} исследуются спектральные свойства матрицы Гессе функции потерь нейронной сети и выводятся архитектурно-зависимые оценки локальной кривизны оптимизационной поверхности. Результаты главы опубликованы в работах~\cite{kiselev2024unraveling,meshkov2024convnets}.

В качестве основного объекта рассматривается матрица Гессе эмпирической функции потерь
\begin{equation}
\mathbf{H}^{(m)}(\mathbf{w})=\nabla_{\mathbf w}^2 \mathcal L_m(\mathbf w).
\end{equation}
Для анализа спектрально значимой части второй производной используется разложение Гаусса--Ньютона
\begin{equation}
\mathbf H(\mathbf w)=\mathbf G(\mathbf w)+\mathbf R(\mathbf w), \qquad
\mathbf G(\mathbf w)=\mathbf J(\mathbf w)^\top \mathbf A(\mathbf w)\mathbf J(\mathbf w),
\end{equation}
где $\mathbf J(\mathbf w)=\partial \mathbf z/\partial \mathbf w$ "--- якобиан выхода модели по параметрам, а $\mathbf A(\mathbf w)=\nabla_{\mathbf z}^2\ell(\mathbf z, \mathbf y)$ "--- матрица кривизны функции потерь по выходу модели.

Для широкого класса матрично-представимых нейронных сетей в работе получена блочная факторизация гаусс--ньютоновской компоненты. Если $\mathbf w=\mathrm{col}(\mathbf w^{(1)}, \dots, \mathbf w^{(L)})$ "--- разбиение параметров по слоям, а
\[
\mathbf S^{(p)}(\mathbf w)=\frac{\partial \mathbf z}{\partial \mathbf z^{(p)}}, \qquad
\mathbf P^{(p)}(\mathbf w)=\frac{\partial \mathbf z^{(p)}}{\partial \mathbf w^{(p)}},
\]
то для $(p, q)$-го блока матрицы $\mathbf G(\mathbf w)$ доказано представление
\begin{equation}
\mathbf G^{(p, q)}(\mathbf w)=
\mathbf P^{(p)}(\mathbf w)^\top
\mathbf S^{(p)}(\mathbf w)^\top
\mathbf A(\mathbf w)
\mathbf S^{(q)}(\mathbf w)
\mathbf P^{(q)}(\mathbf w).
\end{equation}
Полученная факторизация, проиллюстрированная на рисунке~\ref{fig:synopsis-gn-factorization}, позволяет разделить вклад архитектуры сети, функции потерь и параметризации слоев в локальную кривизну.

\begin{figure}[ht]
    \centering
    \resizebox{\linewidth}{!}{\begin{tikzpicture}[
    >=Latex,
    font=\small,
    every node/.style={align=center},
    blockmat/.style={draw, line width=0.5pt},
    flow/.style={-Latex, thick},
    factor/.style={
        draw,
        rounded corners=2pt,
        minimum height=8mm,
        minimum width=14mm,
        inner sep=2pt
    }
]

% -------------------------------------------------
% Left: big block matrix G
% -------------------------------------------------
\def\W{4.8}
\def\H{4.8}

\coordinate (Gsw) at (0,0);
\coordinate (Gne) at (\W,\H);

\draw[blockmat] (Gsw) rectangle (Gne);

% grid 4x4
\foreach \x in {1.2,2.4,3.6}
    \draw[blockmat] (\x,0) -- (\x,\H);
\foreach \y in {1.2,2.4,3.6}
    \draw[blockmat] (0,\y) -- (\W,\y);

% labels inside matrix
\node at (2.4,5.25) {$\mathbf{G}(\mathbf{w})=\big[\mathbf{G}^{(p, q)}(\mathbf{w})\big]_{p, q=1}^{L}$};

\node at (0.6,4.2) {$\mathbf{G}^{(1, 1)}$};
\node at (1.8,4.2) {$\mathbf{G}^{(1, 2)}$};
\node at (3.0,4.2) {$\cdots$};
\node at (4.2,4.2) {$\mathbf{G}^{(1, L)}$};

\node at (0.6,3.0) {$\mathbf{G}^{(2, 1)}$};
\node at (1.8,3.0) {$\mathbf{G}^{(2, 2)}$};
\node at (3.0,3.0) {$\cdots$};
\node at (4.2,3.0) {$\mathbf{G}^{(2, L)}$};

\node at (0.6,1.8) {$\vdots$};
\node at (1.8,1.8) {$\vdots$};
\node at (3.0,1.8) {$\ddots$};
\node at (4.2,1.8) {$\vdots$};

\node at (0.6,0.6) {$\mathbf{G}^{(L, 1)}$};
\node at (1.8,0.6) {$\mathbf{G}^{(L, 2)}$};
\node at (3.0,0.6) {$\cdots$};
\node at (4.2,0.6) {$\mathbf{G}^{(L, L)}$};

% row/col annotations
\node[left=4mm] at (0,2.4) {\rotatebox{90}{\small параметры $p$-го слоя}};
\node[below=4mm] at (2.4,0) {\small параметры $q$-го слоя};

% single arrow from selected block
\coordinate (selected) at (3.0,1.8);

% -------------------------------------------------
% Right: factorization of one block
% -------------------------------------------------
\node[factor] at (8.0,1.8) (Pt) {$\mathbf{P}^{(p)\top}$};
\node[factor, right=0.18cm of Pt] (St) {$\mathbf{S}^{(p)\top}$};
\node[factor, right=0.18cm of St] (A) {$\mathbf{A}(\mathbf{w})$};
\node[factor, right=0.18cm of A] (S) {$\mathbf{S}^{(q)}$};
\node[factor, right=0.18cm of S] (P) {$\mathbf{P}^{(q)}$};

\draw[flow] (selected) -- (Pt.west);

\node[above=6mm of A] {\small факторизация отдельного блока};
\node[below=6mm of A] {$\displaystyle
\mathbf{G}^{(p,q)}(\mathbf{w})
=
\mathbf{P}^{(p)\top}
\mathbf{S}^{(p)\top}
\mathbf{A}(\mathbf{w})
\mathbf{S}^{(q)}
\mathbf{P}^{(q)}
$};

\end{tikzpicture}}
    \caption{Блочная факторизация гаусс--ньютоновской компоненты матрицы Гессе для матрично-представимой нейронной сети}
    \label{fig:synopsis-gn-factorization}
\end{figure}

На основе этого результата для $L$-слойной полносвязной сети получена верхняя оценка спектральной нормы гаусс--ньютоновской компоненты
\begin{equation}
\|\mathbf G(\mathbf w)\|_2
\le
L M_{\mathbf x}^2 M_{\mathbf A}
\left(M_{\mathbf W}M_{\sigma}\right)^{2L-2},
\end{equation}
где $M_{\mathbf x}$ ограничивает норму входа, $M_{\mathbf A}$ "--- спектральную норму матрицы $\mathbf A(\mathbf w)$, $M_{\mathbf W}$ "--- нормы весовых матриц, а $M_{\sigma}$ "--- нормы матриц производных активаций.

Для сверточной нейронной сети доказана аналогичная оценка
\begin{equation}
\|\mathbf G(\mathbf w)\|_2
\le
L M_{\mathbf x}^2 M_{\mathbf A} K_{\max}
\left(M_{\mathcal W}M_{\sigma}\right)^{2L-2},
\end{equation}
где $M_{\mathcal W}$ ограничивает $\ell_1$-нормы сверточных ядер, а множитель $K_{\max}$ определяется максимальным размером рецептивного поля слоя. Тем самым показано, что для сверточных сетей верхняя оценка локальной кривизны зависит от локальности связей и разделения весов.

Кроме того, для полносвязной сети с активацией ReLU (англ. rectified linear unit) и кросс-энтропийной функцией потерь получено частное следствие
\begin{equation}
\|\mathbf G(\mathbf w)\|_2
\le
\sqrt{2}\,L M_{\mathbf x}^2 M_{\mathbf W}^{2L-2}.
\end{equation}

Проведенные вычислительные эксперименты для полносвязных и сверточных архитектур на задаче классификации показывают, что на поздних этапах обучения выполняется соотношение
\[
\|\mathbf G(\mathbf w)\|_2 \approx \|\mathbf H(\mathbf w)\|_2,
\]
что подтверждает возможность использования гаусс--ньютоновской компоненты как практически значимой характеристики локальной кривизны оптимизационной поверхности.

\underline{\textbf{Во второй главе}} исследуется стабилизация ландшафта функции потерь нейронной сети при увеличении объема обучающей выборки. В отличие от классического анализа, основанного на сходимости эмпирического риска как скалярной величины, в главе рассматривается изменение локальной геометрии оптимизационной поверхности в окрестности точки минимума. Результаты главы опубликованы в работах~\cite{kiselev2024unraveling,kiselev2026ispras,kiselev2026conf68mipt,meshkov2024convnets}.

В качестве исходного объекта изучается разность $\mathcal L_{k+1}(\mathbf w)-\mathcal L_k(\mathbf w)$ эмпирических функций потерь на выборках размеров $k$ и $k+1$. В окрестности точки локального минимума $\mathbf w_k^*$ функции $\mathcal L_k(\mathbf w)$ эта разность аппроксимируется формулой второго порядка
\begin{align}
&\mathcal L_{k+1}(\mathbf w)-\mathcal L_k(\mathbf w)
\nonumber \\
&\qquad=
\mathcal L_{k+1}(\mathbf w_k^*)-\mathcal L_k(\mathbf w_k^*)
\nonumber \\
&\qquad+
\mathbf g^{(k+1)}(\mathbf w_k^*)^\top(\mathbf w-\mathbf w_k^*)
\nonumber \\
&\qquad
+
\frac12
(\mathbf w-\mathbf w_k^*)^\top
\bigl(
\mathbf H^{(k+1)}(\mathbf w_k^*)-\mathbf H^{(k)}(\mathbf w_k^*)
\bigr)
(\mathbf w-\mathbf w_k^*).
\end{align}
где $\mathbf g^{(k+1)}(\mathbf w_k^*)$ "--- градиент функции потерь на выборке из $k+1$ объектов, а $\mathbf H^{(k)}(\mathbf w_k^*)$ и $\mathbf H^{(k+1)}(\mathbf w_k^*)$ "--- соответствующие матрицы Гессе. Тем самым изменение ландшафта описывается тремя составляющими: сдвигом значения функции потерь в точке минимума, градиентным рассогласованием (см. иллюстрацию на рисунке~\ref{fig:synopsis-losses-assumption}) и изменением локальной кривизны.

\begin{figure}[ht]
    \centering
    \includegraphics[width=0.62\linewidth]{losses_assumption.pdf}
    \caption{Иллюстрация условия слабого градиентного рассогласования в окрестности локального минимума}
    \label{fig:synopsis-losses-assumption}
\end{figure}

На этой основе вводится общий критерий $p$-сходимости ландшафта функции потерь:
\begin{equation}
\Delta_p(k+1)
=
\int
\left|
\mathcal L_{k+1}(\mathbf w)-\mathcal L_k(\mathbf w)
\right|^p
q(\mathbf w)\,d\mathbf w,
\end{equation}
где функция $q(\mathbf w)$ задает предпочтение по выбору точек в пространстве параметров. Такой критерий позволяет единообразно рассматривать как изменение функции потерь в одной фиксированной точке, так и усредненные характеристики стабилизации поверхности в полном пространстве параметров или в выделенном подпространстве, как это проиллюстрировано на рисунке~\ref{fig:synopsis-landscape-criteria}.

\begin{figure}[ht]
    \centering
    \begin{minipage}{0.32\linewidth}
        \centering
        \includegraphics[width=\linewidth]{criterion_abs.pdf}
        \small а)
    \end{minipage}\hfill
    \begin{minipage}{0.32\linewidth}
        \centering
        \includegraphics[width=\linewidth]{criterion_squared.pdf}
        \small б)
    \end{minipage}\hfill
    \begin{minipage}{0.32\linewidth}
        \centering
        \includegraphics[width=\linewidth]{criterion_squared_subspace.pdf}
        \small в)
    \end{minipage}
    \caption{Иллюстрация введенных критериев стабилизации ландшафта функции потерь: а) абсолютный одноточечный критерий; б) среднеквадратичный критерий в полном пространстве параметров; в) среднеквадратичный критерий в подпространстве наибольшей кривизны}
    \label{fig:synopsis-landscape-criteria}
\end{figure}

Для абсолютного одноточечного критерия
\begin{equation}
\Delta_1(k+1)
=
\left|
\mathcal L_{k+1}(\mathbf w)-\mathcal L_k(\mathbf w)
\right|,
\qquad
\mathbf w\in\mathcal U_R(\mathbf w_k^*),
\end{equation}
получена оценка
\begin{equation}
\Delta_1(k+1)
\le
\frac{2M_\ell+M_{\mathbf g}R+M_{\mathbf H}R^2}{k+1}
=
\mathcal O(k^{-1}),
\end{equation}
где $M_\ell$, $M_{\mathbf g}$ и $M_{\mathbf H}$ ограничивают соответственно значения функции потерь на одном объекте, градиента и матрицы Гессе в окрестности минимума.

Для среднеквадратичного критерия
\begin{equation}
\Delta_2(k+1)
=
\int
\big(
\mathcal L_{k+1}(\mathbf w)-\mathcal L_k(\mathbf w)
\big)^2
q(\mathbf w)\,d\mathbf w
\end{equation}
при выборе гауссовского распределения
\begin{equation}
q(\mathbf w)=\mathcal N(\mathbf w_k^*, \sigma^2\mathbf I_N)
\end{equation}
доказана оценка
\begin{equation}
\Delta_2(k+1)
\le
\frac{
12M_\ell^2
+
3\sigma^2M_{\mathbf g}^2
+
3\sigma^4(N^2+2N)M_{\mathbf H}^2
}{(k+1)^2}
=
\mathcal O(k^{-2}).
\end{equation}
Таким образом, после усреднения квадрата изменения ландшафта по распределению точек получается более быстрая скорость убывания по сравнению с одноточечным критерием. Множитель $(k+1)^{-2}$ возникает из тождества усреднения до квадратичной аппроксимации; матрица Гессе ограничивает числитель.

Далее вводится подпространство наибольшей кривизны
\begin{equation}
\mathcal S_D=\operatorname{Im}(\mathbf U_D),
\qquad
\mathbf U_D=[\mathbf u_1, \dots, \mathbf u_D],
\end{equation}
натянутое на собственные векторы $\mathbf u_1, \dots, \mathbf u_D$, соответствующие $D$ наибольшим собственным значениям матрицы Гессе $\mathbf H^{(k)}(\mathbf w_k^*)$. В этом подпространстве рассматривается критерий
\begin{equation}
\Delta_2^{(D)}(k+1)
=
\int
\big(
\mathcal L_{k+1}(\mathbf w)-\mathcal L_k(\mathbf w)
\big)^2
q(\mathbf w)\,d\mathbf w,
\qquad
\mathbf w\in \mathbf w_k^*+\mathcal S_D.
\end{equation}
При гауссовском распределении в координатах подпространства доказано, что
\begin{equation}
\Delta_2^{(D)}(k+1)
\le
\frac{
12M_\ell^2
+
3\sigma^2M_{\mathbf g}^2
+
3\sigma^4(D^2+2D)M_{\mathbf H}^2
}{(k+1)^2}
=
\mathcal O(k^{-2}).
\end{equation}
Тем самым показано, что ограничение анализа подпространством главных собственных направлений матрицы Гессе сохраняет тот же порядок сходимости, что и в полном пространстве параметров, но позволяет заменить размерность $N$ на существенно меньшую размерность $D$.

В частном случае отсутствия нулевого и линейного вкладов критерий в подпространстве наибольшей кривизны выражается непосредственно через приращения ведущих собственных значений матрицы Гессе:
\begin{equation}
\Delta_{2}^{(D)}(k+1)
=
\frac{\sigma^4}{4}
\left(
2\sum_{i=1}^{D}
\left(
\lambda_i^{(k+1)}-\lambda_i^{(k)}
\right)^2
+
\left(
\sum_{i=1}^{D}
\left(
\lambda_i^{(k+1)}-\lambda_i^{(k)}
\right)
\right)^2
\right).
\end{equation}
Это дает спектральную интерпретацию стабилизации ландшафта функции потерь.

Вычислительный эксперимент для полносвязной нейронной сети на задаче классификации MNIST проверяет порядок. Среднеквадратичный критерий в полном пространстве и подпространственные критерии убывают близко к $\mathcal O(k^{-2})$. Этот порядок следует из тождества усреднения; эксперимент не измеряет константу, зависящую от матрицы Гессе. Абсолютный одноточечный критерий также убывает, но его степенная оценка неустойчива. Спектр матрицы Гессе отделяет несколько старших собственных значений от более пологого хвоста.

\underline{\textbf{В третьей главе}} результаты, полученные ранее для спектральных свойств матрицы Гессе и для критериев стабилизации ландшафта функции потерь, объединяются в единую схему, позволяющую построить архитектурно-зависимые критерии достаточного размера выборки для параметрических моделей машинного обучения. Результаты главы опубликованы в работах~\cite{kiselev2025conf67mipt,kiselev2025mmro}.

Исходным пунктом является разложение пообъектной матрицы Гессе в форме Гаусса--Ньютона
\begin{equation}
\mathbf H_i(\mathbf w)=\mathbf G_i(\mathbf w)+\mathbf R_i(\mathbf w),
\end{equation}
где $\mathbf G_i(\mathbf w)$ есть гаусс--ньютоновская компонента, а $\mathbf R_i(\mathbf w)$ "--- остаточный член. Для фиксированной архитектуры $\mathcal A$ вводится локальная архитектурно-зависимая константа кривизны
\begin{equation}
M_{\mathbf G}(\mathcal A)
=
\sup_{\mathbf w\in\mathcal U_R(\mathbf w^*)}
\sup_i
\|\mathbf G_i(\mathbf w)\|_2,
\end{equation}
а также величина
\begin{equation}
\delta_R
=
\sup_{\mathbf w\in\mathcal U_R(\mathbf w^*)}
\sup_i
\|\mathbf R_i(\mathbf w)\|_2,
\end{equation}
характеризующая вклад остаточного члена. Тогда для любой точки локальной окрестности минимума справедливо неравенство
\begin{equation}
\|\mathbf H_i(\mathbf w)\|_2
\le
M_{\mathbf G}(\mathcal A)+\delta_R.
\end{equation}
Это позволяет во всех критериях стабилизации ландшафта заменить общую константу $M_{\mathbf H}$ на величину $M_{\mathbf G}(\mathcal A)+\delta_R$, зависящую от архитектуры модели.

На этой основе достаточный размер выборки определяется через выбранный критерий стабилизации ландшафта функции потерь:
\begin{equation}
m^*
=
\inf\{m\in\mathbb N:\Delta(m)\le\varepsilon\},
\end{equation}
где $\varepsilon>0$ "--- заранее заданный уровень точности. Тем самым выборка считается достаточной, если при дальнейшем увеличении числа обучающих объектов изменение локальной геометрии функции потерь становится несущественным по выбранному критерию.

Для абсолютного одноточечного критерия получена оценка достаточного размера выборки вида
\begin{equation}
m_1^*(\varepsilon, \mathcal A)
=
\frac{
2M_{\ell}
+
M_{\mathbf g}R
+
\left(M_{\mathbf G}(\mathcal A)+\delta_R\right)R^2
}{\varepsilon},
\end{equation}
где $M_{\ell}$ ограничивает значения функции потерь на отдельных объектах, $M_{\mathbf g}$ "--- величину градиентного рассогласования, а $R$ "--- радиус локальной окрестности минимума.

Для среднеквадратичного критерия в полном пространстве параметров при гауссовской функции предпочтения
\begin{equation}
q(\mathbf w)=\mathcal N(\mathbf w^*, \sigma^2\mathbf I_N)
\end{equation}
получена оценка
\begin{equation}
m_2^*(\varepsilon, \mathcal A)
=
\left[
\frac{
12M_{\ell}^2
+
3\sigma^2M_{\mathbf g}^2
+
3\sigma^4(N^2+2N)\left(M_{\mathbf G}(\mathcal A)+\delta_R\right)^2
}{\varepsilon}
\right]^{1/2}.
\end{equation}
Аналогично, для среднеквадратичного критерия в подпространстве наибольшей кривизны размерности $D$ доказана оценка
\begin{equation}
m_3^*(\varepsilon, \mathcal A, D)
=
\left[
\frac{
12M_{\ell}^2
+
3\sigma^2M_{\mathbf g}^2
+
3\sigma^4(D^2+2D)\left(M_{\mathbf G}(\mathcal A)+\delta_R\right)^2
}{\varepsilon}
\right]^{1/2}.
\end{equation}
По сравнению с полным пространством параметров здесь размерностный множитель зависит не от полной размерности модели $N$, а только от размерности выбранного информативного подпространства $D$.

Подстановка полученных в первой главе архитектурно-зависимых оценок для $M_{\mathbf G}(\mathcal A)$ дает явные критерии достаточного размера выборки для полносвязных и сверточных нейронных сетей. Для $L$-слойной полносвязной сети без смещений локальная константа кривизны имеет вид
\begin{equation}
M_{\mathbf G}^{\mathrm{FC}}
=
L M_{\mathbf x}^2 M_{\mathbf A}(M_{\mathbf W}M_{\sigma})^{2L-2},
\end{equation}
а для $L$-слойной сверточной сети без смещений
\begin{equation}
M_{\mathbf G}^{\mathrm{CNN}}
=
L M_{\mathbf x}^2 M_{\mathbf A}K_{\max}(M_{\mathcal W}M_{\sigma})^{2L-2}.
\end{equation}
Тем самым достаточный размер выборки выражается через глубину архитектуры, нормы весовых операторов, свойства функций активации и, в сверточном случае, размер рецептивного поля.

Полученные формулы показывают, что при прочих равных увеличение глубины приводит к росту необходимого объема данных для стабилизации ландшафта функции потерь. В полносвязных сетях этот рост связан с накоплением кривизны в длинных произведениях плотных операторов, тогда как в сверточных сетях существенную роль играет локальность связей, отражаемая множителем $K_{\max}$. В результате архитектура модели интерпретируется не только как класс реализуемых функций, но и как фактор, определяющий скорость стабилизации локальной геометрии задачи оптимизации.

Помимо строгих критериев достаточного размера выборки, в главе предложена более общая интерпретация полученных результатов в терминах законов масштабирования параметрических моделей. Для этого рассматривается декомпозиция среднего значения популяционной функции потерь в обученной точке $\mathbf w_m^*$:
\begin{align}
\mathbb E_{\mathfrak D_m}\mathcal L(\mathbf w_m^*)
&=
\underbrace{\mathcal L(\mathbf w_{\infty}^*)}_{E}
+
\underbrace{\Phi(N)}_{\text{ошибка выразимости}}
+
\underbrace{\Bigl(\mathbb E_{\mathfrak D_m}\mathcal L(\mathbf w_m^*)-\inf_{\mathbf w\in\mathcal W_N}\mathcal L(\mathbf w)\Bigr)}_{\mathcal E(m)},
\\
\Phi(N)
&=
\inf_{\mathbf w\in\mathcal W_N}\mathcal L(\mathbf w)-\mathcal L(\mathbf w_{\infty}^*).
\end{align}
Здесь $E$ есть неустранимая ошибка задачи, $\Phi(N)$ "--- ошибка выразимости класса моделей с $N$ параметрами, а $\mathcal E(m)\geqslant 0$ "--- ошибка оценивания, то есть та часть ошибки, которую может устранить увеличение объема данных. Структурное предположение о локальном статистическом режиме имеет вид
\begin{equation}
\mathcal E(m)
\le
\frac{C\left(M_{\mathbf G}(\mathcal A)+\delta_R\right)}{m^{\rho}},
\qquad
\rho\in(0, 1],
\end{equation}
и приводит к схеме локального закона масштабирования
\begin{equation}
\mathbb E_{\mathfrak D_m}\mathcal L(\mathbf w_m^*)
\le
E
+
\Phi(N)
+
\frac{C\left(M_{\mathbf G}(\mathcal A)+\delta_R\right)}{m^{\rho}}.
\end{equation}

Для линейной регрессии с гауссовским распределением признаков предположение выполняется точно: средний избыточный риск оценки наименьших квадратов равен
\begin{equation}
\mathcal E(m)=\frac{\sigma^2 n}{2(m-n-1)}=\frac{\operatorname{tr}(\mathbf H^{-1}\mathbf C)}{2m}\left(1+\frac{n+1}{m-n-1}\right),
\end{equation}
то есть $\rho=1$, а константа определяется отношением шума градиентов к кривизне $\operatorname{tr}(\mathbf H^{-1}\mathbf C)$. Архитектурно-зависимая константа входит в верхнюю оценку этой величины: $\operatorname{tr}(\mathbf H^{-1}\mathbf C)\leqslant \|\mathbf C\|_2\operatorname{tr}(\mathbf H^{-1})\leqslant a_{\min}^{-1}\,\mathbb E_{(\mathbf x, \mathbf y)}\|\mathbf s\|^2 M_{\mathbf G}(\mathcal A)\operatorname{tr}(\mathbf H^{-1})$.

Для линейной модели с неполным набором признаков, использующей первые $N$ собственных направлений матрицы Гессе $\boldsymbol\Lambda=\operatorname{diag}(\lambda_j)$, схема реализуется точно со всеми тремя слагаемыми:
\begin{equation}
\mathbb E_{\mathfrak D_m}\,\mathcal L(\hat{\mathbf w}_{N, m})
=
\frac{\sigma^2}{2}
+
\Phi(N)
+
\frac{\bigl(\sigma^2+2\Phi(N)\bigr)N}{2(m-N-1)},
\qquad
\Phi(N)=\frac12\sum_{j>N}\lambda_j(\theta_j^\star)^2.
\end{equation}
Если энергия истинной зависимости убывает по собственным направлениям кривизны как $\lambda_j(\theta_j^\star)^2\asymp j^{-b}$, $b>1$, то показатель ошибки выразимости равен $\alpha=b-1$, при оптимальном размере модели риск убывает по объему данных как $m^{-(b-1)/b}$, а при вычислительном бюджете $\mathcal C\propto Nm$ оптимальный размер модели растет как $\mathcal C^{1/(\alpha+2)}$.

В общем случае при вычислительном бюджете $\mathcal C\propto Nm$ и росте коэффициента ошибки оценивания с размером модели как $K(N)\sim BN^{\gamma}$ оптимальное соотношение между размером модели и размером выборки удовлетворяет асимптотикам
\begin{equation}
N^*
\propto
\mathcal C^{\frac{\rho}{\alpha+\gamma+\rho}},
\qquad
m^*
\propto
\mathcal C^{\frac{\alpha+\gamma}{\alpha+\gamma+\rho}}.
\end{equation}
В точно решаемой модели $\gamma=1$: показатель описывает рост с размером модели отношения шума к кривизне, то есть эффективного числа параметров, а не спектральной нормы кривизны, которая от $N$ не зависит.

Проведенные вычислительные эксперименты на полносвязных и сверточных архитектурах для задачи классификации MNIST согласуются с теми выводами главы, которые задают порядок или монотонность. Во-первых, критерий стабилизации ландшафта $\Delta_2(m)$ убывает с ростом объема выборки с порядком, близким к $\mathcal O(m^{-2})$. Этот порядок следует из усреднения по объектам. Во-вторых, внутри каждого семейства архитектур достаточный размер выборки возрастает с увеличением глубины. В-третьих, в локальном режиме линейной и логистической регрессии константа избыточного риска, вычисленная по спектру, согласуется с методом Монте--Карло, и показатель $k^{-1}$ по кривой не подбирается.

\underline{\textbf{В четвертой главе}} рассматриваются критерии достаточного размера выборки для вероятностных линейных моделей. В отличие от предыдущих глав, где достаточность выборки исследовалась через стабилизацию ландшафта функции потерь и его локальной геометрии, здесь изучается вероятностная постановка, в которой основным объектом анализа становится стабилизация оценок параметров и соответствующих апостериорных распределений. Результаты главы опубликованы в работах~\cite{kiselev2025ssdlb,kiselev2025ssdpdp,kiselev2024conf66mipt}.

Рассматривается вероятностная модель линейной регрессии
\begin{equation}
p(y | \mathbf x, \mathbf w)=\mathcal N(y | \mathbf w^\top \mathbf x, \sigma^2),
\end{equation}
где $\mathbf x\in\mathbb R^n$ "--- вектор признаков, $y\in\mathbb R$ "--- отклик, $\mathbf w\in\mathbb R^n$ "--- вектор параметров, а $\sigma^2>0$ "--- дисперсия шума. Для выборки $\mathfrak D_k=(\mathbf X_k, \mathbf y_k)$ оценка максимального правдоподобия определяется как решение задачи наименьших квадратов
\begin{equation}
\hat{\mathbf w}_k
=
\argmin_{\mathbf w\in\mathbb R^n}
\|\mathbf y_k-\mathbf X_k\mathbf w\|_2^2.
\end{equation}
При гауссовском априорном распределении параметров апостериорное распределение имеет сопряженный вид
\begin{equation}
p(\mathbf w | \mathfrak D_k)=\mathcal N(\mathbf w | \mathbf m_k, \boldsymbol\Sigma_k),
\end{equation}
где параметры $\mathbf m_k$ и $\boldsymbol\Sigma_k$ явно выражаются через $\mathbf X_k$ и $\mathbf y_k$.

На этой основе в главе вводятся четыре критерия достаточного размера выборки, основанные на различных характеристиках стабилизации решения. Первые два критерия, проиллюстрированные на рисунке~\ref{fig:likelihood}, используют поведение функции правдоподобия в точке оценки максимального правдоподобия. D-критерий основан на уменьшении дисперсии правдоподобия при варьировании подвыборок, а M-критерий "--- на уменьшении изменения его математического ожидания при переходе от выборки размера $k$ к выборке размера $k+1$.

\begin{figure}[ht]
    \centering
    \includegraphics[width=0.5\linewidth]{likelihood.pdf}
    \caption{Визуализация критериев D-достаточности и M-достаточности размера выборки в вероятностной модели линейной регрессии}
    \label{fig:likelihood}
\end{figure}

Оставшиеся два критерия, проиллюстрированные на рисунке~\ref{fig:posterior}, используют близость апостериорных распределений параметров на схожих подвыборках. KL-критерий определяется через дивергенцию Кульбака--Лейблера
\begin{equation}
D_{\mathrm{KL}}(p_k\|p_{k+1})
=
\int p_k(\mathbf w)\log\frac{p_k(\mathbf w)}{p_{k+1}(\mathbf w)}\,d\mathbf w,
\end{equation}
а S-критерий "--- через функцию близости (англ. s-score), нормированную так, что для совпадающих распределений она равна единице. Тем самым достаточный размер выборки определяется через стабилизацию либо правдоподобия, либо полного апостериорного распределения параметров.

\begin{figure}[ht]
    \centering
    \includegraphics[width=0.88\linewidth]{posterior_ru.pdf}
    \caption{Визуализация критериев KL-достаточности и S-достаточности размера выборки в вероятностной модели линейной регрессии}
    \label{fig:posterior}
\end{figure}

В главе формализуется понятие корректности критерия достаточности: критерий называется корректным, если соответствующая функция сходится к своему предельному значению при увеличении объема выборки. Для M-, KL- и S-критериев доказываются теоремы корректности, связывающие ее со сходимостью моментов апостериорного распределения. В частности, показывается, что если
\begin{equation}
\|\mathbf m_{k+1}-\mathbf m_k\|_2\to 0,
\qquad
\|\boldsymbol\Sigma_{k+1}-\boldsymbol\Sigma_k\|_F\to 0
\quad\text{при}\quad k\to\infty,
\end{equation}
то M- и KL-критерии являются корректными, а для S-критерия достаточно сходимости апостериорных средних.

Далее для вероятностной линейной регрессии с гауссовским априором доказывается теорема, дающая достаточные условия для выполнения этих сходимостей. Предполагается ограниченность признаков и отклика, а также условие роста минимального собственного значения матрицы Грама:
\begin{equation}
\lambda_{\min}(\mathbf X_k^\top \mathbf X_k)=\omega(\sqrt{k}).
\end{equation}
При этих предположениях устанавливается, что
\begin{equation}
\|\mathbf m_{k+1}-\mathbf m_k\|_2\to 0,
\qquad
\|\boldsymbol\Sigma_{k+1}-\boldsymbol\Sigma_k\|_F\to 0
\quad\text{при}\quad k\to\infty.
\end{equation}
Тем самым получаются проверяемые условия, обеспечивающие корректность введенных критериев достаточности в модели линейной регрессии.

Помимо стабилизации самих распределений параметров, в главе исследуется стабилизация полного байесовского прогноза на отложенной тестовой выборке. Пусть
\begin{equation}
p(\mathbf y_{\mathrm{test}} | \mathbf X_{\mathrm{test}}, \mathbf X_k, \mathbf y_k)
\end{equation}
обозначает предиктивное распределение, полученное по обучающей подвыборке размера $k$. Доказывается, что при тех же условиях на данные и матрицу признаков
\begin{equation}
D_{\mathrm{KL}}
\left(
p(\mathbf y_{\mathrm{test}} | \mathbf X_{\mathrm{test}}, \mathbf X_k, \mathbf y_k)
\big\|
p(\mathbf y_{\mathrm{test}} | \mathbf X_{\mathrm{test}}, \mathbf X_{k+1}, \mathbf y_{k+1})
\right)
\to 0
\quad\text{при}\quad k\to\infty.
\end{equation}
Это означает, что увеличение объема обучающей выборки приводит не только к стабилизации апостериорного распределения параметров, но и к стабилизации байесовского прогноза на новых данных.

Вычислительный эксперимент проведен на синтетической выборке линейной регрессии, на метках $\pm 1$, порожденных логистической моделью и оцененных по гауссовскому правдоподобию в точке оценки наименьших квадратов, и на одном наборе данных Liver Disorders. При увеличении размера подвыборки функции $D(k)$, $M(k)$ и $KL(k)$ стремятся к нулю, а функция $S(k)$ стремится к единице. Это согласуется с корректностью критериев и не задает скорость убывания: теоремы гарантируют существование достаточного размера выборки, а не порядок.

Таким образом, в содержательной части автореферата последовательно раскрывается общая логика диссертационного исследования. Сначала формулируется задача определения достаточного размера обучающей выборки для параметрических моделей, затем исследуются архитектурно-зависимые свойства локальной кривизны оптимизационной поверхности, после чего строятся критерии стабилизации ландшафта функции потерь и выводятся критерии достаточного размера выборки. Отдельно рассматривается вероятностная линейная регрессия как частный случай параметрических моделей, для которого достаточность выборки описывается через стабилизацию распределения оптимальных параметров и байесовского прогноза.

\FloatBarrier
\pdfbookmark{Заключение}{conclusion}                                  % Закладка pdf

В \underline{\textbf{заключении}} приведены основные результаты работы, которые заключаются в следующем:
%% Согласно ГОСТ Р 7.0.11-2011:
%% 5.3.3 В заключении диссертации излагают итоги выполненного исследования, рекомендации, перспективы дальнейшей разработки темы.
%% 9.2.3 В заключении автореферата диссертации излагают итоги данного исследования, рекомендации и перспективы дальнейшей разработки темы.
%% Формулировки согласованы с заключением диссертации.
\begin{enumerate}
    \item Для нейронных сетей получена блочная факторизация гаусс"--~ньютоновской компоненты матрицы Гессе в классе матрично-представимых моделей и выведены архитектурно-зависимые верхние оценки ее спектральной нормы для полносвязных и сверточных архитектур.
    \item Введен общий критерий $p$-сходимости ландшафта функции потерь и его частные случаи: абсолютный одноточечный критерий, среднеквадратичный критерий в полном пространстве параметров и среднеквадратичный критерий в подпространстве наибольшей кривизны.
    \item Порядок убывания этих критериев следует из тождества усреднения: одноточечный критерий имеет порядок $\mathcal{O}(k^{-1})$, среднеквадратичный критерий и его подпространственный аналог имеют порядок $\mathcal{O}(k^{-2})$. Квадратичная аппроксимация и матрица Гессе входят в константу, а не в порядок.
    \item Построены критерии достаточного размера выборки на основе стабилизации ландшафта функции потерь. Верхние оценки локальной кривизны через спектральную норму гаусс"--~ньютоновской компоненты дают архитектурно-зависимые оценки достаточного объема данных для полносвязных и сверточных моделей.
    \item Предложена схема законов масштабирования, в которой ошибка оценивания определяется отношением шума градиентов к кривизне оптимизационной поверхности. Для линейной регрессии схема выполняется точно; для линейной модели с неполным набором признаков показатели по размеру модели и объему данных выражены через спектр матрицы Гессе.
    \item Для вероятностных линейных моделей построены критерии достаточного размера выборки по дисперсии и математическому ожиданию функции правдоподобия, дивергенции Кульбака"--~Лейблера между апостериорными распределениями и функции близости (англ. s-score). Получены условия их корректности.
    \item Вычислительные эксперименты согласуются с теорией там, где она дает проверяемый порядок или точную формулу. Наклон среднеквадратичного критерия близок к $-2$ и проверяет порядок из тождества усреднения. Для линейной модели избыточный риск совпадает с точной формулой, а вероятностные критерии сходятся к своим пределам. Теоремы о существовании достаточного размера выборки не задают скорость его убывания.
\end{enumerate}


\pdfbookmark{Литература}{bibliography}                                % Закладка pdf

\ifdefmacro{\microtypesetup}{\microtypesetup{protrusion=false}}{} % не рекомендуется применять пакет микротипографики к автоматически генерируемому списку литературы
\urlstyle{rm}                               % ссылки URL обычным шрифтом
\ifnumequal{\value{bibliosel}}{0}{% Встроенная реализация с загрузкой файла через движок bibtex8
    \renewcommand{\bibname}{\large \bibtitleauthor}
    \nocite{*}
    \insertbiblioauthor           % Подключаем Bib-базы
    %\insertbiblioexternal   % !!! bibtex не умеет работать с несколькими библиографиями !!!
}{% Реализация пакетом biblatex через движок biber
    % Цитирования.
    %  * Порядок перечисления определяет порядок в библиографии (только внутри подраздела, если `\insertbiblioauthorgrouped`).
    %  * Если не соблюдать порядок "как для \printbibliography", нумерация в `\insertbiblioauthor` будет кривой.
    %  * Если цитировать каждый источник отдельной командой "--- найти некоторые ошибки будет проще.
    %
    % \nocite{}
    \nocite{kiselev2026ispras}
    \nocite{kiselev2025ssdpdp}
    \nocite{kiselev2025ssdlb}
    \nocite{kiselev2024unraveling}
    \nocite{kiselev2026conf68mipt}
    \nocite{kiselev2025mmro}
    \nocite{kiselev2025conf67mipt}
    \nocite{meshkov2024convnets}
    \nocite{kiselev2024conf66mipt}

    % Автореферат содержит только работы автора; ссылок на чужие работы и общего списка литературы нет
    \begin{refcontext}[labelprefix={}]
        \ifnum \value{bibgrouped}>0
            \insertbiblioauthorgrouped
        \else
            \insertbiblioauthor
        \fi
    \end{refcontext}
}
\ifdefmacro{\microtypesetup}{\microtypesetup{protrusion=true}}{}
\urlstyle{tt}                               % возвращаем установки шрифта ссылок URL
